\newif\ifglobal

%%%%%%%%%%%%%%%%%%%%%%%
%%% Compile to view %%%
%%%%%%%%%%%%%%%%%%%%%%%

%%%%%%%%%%%%%%%%%%%%%%%%%%%%%%%%%%%%%%%%%%%%%%%%%%%
%%% Readme:                                     %%%
%%% https://www.overleaf.com/read/bwdjvfjksvjy/ %%%
%%%%%%%%%%%%%%%%%%%%%%%%%%%%%%%%%%%%%%%%%%%%%%%%%%%

% >>> M?: marks a module setting number or important notice

% >>> M4: Set {./relative-path} to external files for this module <<<
\def \modulefiles {./05-IOMUXGPIO}
% To add an external file, use:
% (Replace '---' with actual filename)
% '\input{\modulefiles/tables/---.tex}' for tables
% '\includegraphics{\modulefiles/figures/---}' for figures
% '\subfile{\modulefiles/---}' for register files


% >>> M1: This module is in global project? <<<
% 'YES' - uncomment; 'NO' - comment out
\globaltrue

\ifglobal
    \documentclass[main\_\_CN.tex]{subfiles}

\else
    % Font "FandolSong-Regular" used by ctex does not have GSUB table.
    % It causes warnings that do not affect compilation and can be ignored.
    \PassOptionsToPackage{quiet}{fontspec}

    \documentclass[openany, 10pt]{book}

    % >>> M2: In ./00-shared/preamble-trm-module, also find
    %         settings for visibility of todo notes, watermark <<<
    \usepackage{preamble-trm-module}


    % Enable Chinese typesetting
    \usepackage{ctex}

    % >>> M3: Temporary labels in Chinese
    %         you can add your own labels there <<<
    \myexternaldocument{./00-shared/temp-labels-cn}

\fi %\ifglobal


\setcounter{secnumdepth}{4}


% >>> M5: If this module needs any updates in
% './00-shared/preamble-trm-module.sty'
% (include additional package, etc.), contact your module PM
% or create the file './00-shared/changelog.tex' make the
% necessary updates yourself and document them in this file <<<

% Make text left-aligned and allow latex to break pages naturally, comment out for full justification
\raggedright
\raggedbottom

\begin{document}


% Sets language into which pre-defined bits of text are translated
\selectlanguage{chinese}
\setchineseformatting

% Do not delete this block
\ifglobal
% Add the following front matter if
% this module is outside of global project
\else
    \InputIfFileExists{readme}{}{}
    \listoftodos
    {\let\clearpage\relax \listofdonetodos}
    \tableofcontents
    %\listoftables
    %\listoffigures
\fi



%%%%%%%%%%%%%%%%%%%%%%%%%
%%% TRM Module Begins %%%
%%%%%%%%%%%%%%%%%%%%%%%%%

\hypertarget{iomuxgpio}{}
\chapter{IO MUX 和 GPIO 交换矩阵}
\label{mod:iomuxgpio}

\section{概述}

ESP32-S2 芯片有 43 个物理 GPIO pad。每个 pad 都可用作一个通用 IO，或连接一个内部的外设信号。IO MUX、RTC IO MUX 和 GPIO 交换矩阵用于将信号从外设传输至 GPIO pad。这些模块共同组成了芯片的 IO 控制。

{\bfseries 注意：这 43 个物理 GPIO pad 的序列号为：0 \verb+~+ 21, 26 \verb+~+ 46。其中 GPIO46 仅用作输入管脚，其他的既可以作为输入又可以作为输出管脚。}

本章节主要描述了如何通过 GPIO 交换矩阵、IO MUX 或 RTC IO MUX 实现 43 个数字 pad（控制信号：DRV、IE、OE、WPU、WPD 等）与内部信号的选择和连接，这些内部信号包括：

\begin{itemize}
    \item 116 个数字外设输入信号（控制信号：SIG\_IN\_SEL、SIG\_OUT\_SEL、IE、OE 等）
    \item 182 个数字外设输出信号（控制信号：SIG\_IN\_SEL、SIG\_OUT\_SEL、IE、OE 等）
    \item 快速外设输入和输出信号（控制信号：IE、OE 等）
    \item 22 个 RTC GPIO 信号
\end{itemize}

\begin{figure}[H]
    \centering
    \includegraphics[width=0.76\textwidth]{\modulefiles/figures/iomux_gpio_matrix_all}
    \caption{IO MUX、RTC IO MUX 和 GPIO 交换矩阵结构框图}
    \label{Figure:iomux_gpio_matrix_all}
\end{figure}

图 \ref{Figure:iomux_gpio_matrix_all}为IO MUX、RTC IO MUX 和 GPIO 交换矩阵结构框图。

\begin{enumerate}

\item IO MUX 中每个 GPIO pad 有一个寄存器 \hyperref[regdesc:IOMUXnREG]{IO\_MUX\_\regindex{n}\_REG}，每个 pad 可以配置成：
\begin{itemize}
    \item GPIO 功能，连接 GPIO 交换矩阵；
    \item 直连功能，旁路 GPIO 交换矩阵。
\end{itemize}
快速信号如 SPI、JTAG、UART 等可以旁路 GPIO 交换矩阵以实现更好的高频数字特性。所以高速信号会直接通过 IO MUX 输入和输出。

章节 \ref{sec:iomux_pads} 列出了所有 GPIO pad 的 IO MUX 功能。

\item GPIO 交换矩阵是外设输入输出信号和管脚之间的全交换矩阵。
\begin{itemize}
\item 芯片输入方向：116 个外设输入信号都可以选择任意一个 GPIO pad 的输入信号。
\item 芯片输出方向：每个 GPIO pad 的输出信号可来自 182 个外设输出信号中的任意一个。
\end{itemize}

章节 \ref{sec:peripheral_signals} 列出了 GPIO 交换矩阵的外设信号。

\item RTC IO MUX 用于控制 GPIO pad 的低功耗和模拟功能。只有部分 GPIO pad 具有这些功能。

章节 \ref{sec:rtcmux_pins} 列出了 RTC IO MUX 功能。
\end{enumerate}

\section{通过 GPIO 交换矩阵的外设输入}


\subsection{概述}

为实现通过 GPIO 交换矩阵接收外设输入信号，需要配置 GPIO 交换矩阵从 43 个 GPIO（0 \verb+~+ 21, 26 \verb+~+ 46）中获取外设输入信号的索引号，见交换矩阵表格 \ref{table:peripheral_signals}。并需要配置外设输入选择通过 GPIO 交换矩阵接收输入信号。


\subsection{信号同步}
\label{sec:inputconfig}

如图 \ref{Figure:iomux_gpio_matrix_all} 所示，对于信号输入，外部输入信号从 pad 输入，经 GPIO SYNC 模块同步至 APB 总线时钟后进入 GPIO 交换矩阵。外部输入信号也可以通过 IO MUX 直接进入外设，但信号无法经由 GPIO SYNC 模块同步。

\begin{figure}[h]
    \centering
    \fbox{\includegraphics[width=0.69\textwidth]{\modulefiles/figures/gpio_input_sync}}
    \caption{GPIO 输入经 clock 上升沿或下降沿同步}
    \label{Figure:gpio_input_sync}
\end{figure}

\newpage

GPIO SYNC 模块的功能如图 \ref{Figure:gpio_input_sync} 所示。其中，negative sync 为 GPIO input 经过 APB clock 的下降沿同步，positive sync 为 GPIO input 经过 APB clock 上升沿同步。


\subsection{功能描述}

把某个外设输入信号 \regindex{Y} 绑定到某个 GPIO pad \regindex{X} 的配置过程为：

\begin{enumerate}
    \item 在 GPIO 交换矩阵中配置外设信号 \regindex{Y} 的 \hyperref[regdesc:GPIOFUNCNINSELCFGREG]{GPIO\_FUNC\regindex{y}\_IN\_SEL\_CFG\_REG} 寄存器：
    \begin{itemize}
        \item 置位 \hyperref[fielddesc:GPIOSIGNINSEL]{GPIO\_SIG\regindex{y}\_IN\_SEL} 选择通过 GPIO 交换矩阵接收外部输入信号；
    \end{itemize}
    \begin{itemize}
        \item 设置 \hyperref[fielddesc:GPIOFUNCNINSEL]{GPIO\_FUNC\regindex{y}\_IN\_SEL} 为要读取的 GPIO pad \regindex{X} 的值。
    \end{itemize}
    {\bfseries 注意: } 并不是所有外设信号都有有效的 \hyperref[fielddesc:GPIOSIGNINSEL]{GPIO\_SIG\regindex{y}\_IN\_SEL}，即有些外设信号没有图 \ref{Figure:iomux_gpio_matrix_all} 中所示的 MUX（见图下方的注释 \textcircled{1}）。这些外设信号只能通过 GPIO 交换矩阵接收外部输入信号。
        \item 置位寄存器 \hyperref[fielddesc:IOMUXFILTEREN]{IO\_MUX\_FILTER\_EN} 使能 pad 输入信号滤波，如图 \ref{Figure:gpio_filter} 所示。只有当输入信号的有效宽度大于两个时钟周期时，输入信号才会被采样。否则，输入信号将会被滤掉。

\begin{figure}[h!]
    \centering
    \fbox{\includegraphics[width=0.9\textwidth]{\modulefiles/figures/gpio_filter}}
    \caption{GPIO 输入信号滤波时序图}
    \label{Figure:gpio_filter}
\end{figure}

        \item 配置 GPIO pad \regindex{X} 的 \hyperref[regdesc:GPIOPINNREG]{GPIO\_PIN\regindex{x}\_REG} 来同步 GPIO input，过程如下：
    \begin{itemize}
        \item 如图 \ref{Figure:gpio_input_sync} 所示, 配置 \hyperref[fielddesc:GPIOPINNSYNC1BYPASS]{GPIO\_PIN\regindex{x}\_SYNC1\_BYPASS} 使能输入信号第一拍为上升沿或者下降沿同步。
    \end{itemize}
    \begin{itemize}
        \item 如图 \ref{Figure:gpio_input_sync} 所示, 配置 \hyperref[fielddesc:GPIOPINNSYNC2BYPASS]{GPIO\_PIN\regindex{x}\_SYNC2\_BYPASS} 使能输入信号第二拍为上升沿或者下降沿同步。
    \end{itemize}
    \item 配置 GPIO pad \regindex{X} 的 \hyperref[regdesc:IOMUXnREG]{IO\_MUX\_\regindex{x}\_REG} 来使能 pad 输入功能，过程如下：
    \begin{itemize}
        \item 置位 \hyperref[fielddesc:IOMUXFUNIE]{IO\_MUX\_FUN\_IE} 使能输入；
        \item 置位或清零 \hyperref[fielddesc:IOMUXFUNWPU]{IO\_MUX\_FUN\_WPU} 和 \hyperref[fielddesc:IOMUXFUNWPD]{IO\_MUX\_FUN\_WPD}，使能或关闭内部上拉/下拉电阻器。
    \end{itemize}
\end{enumerate}

例如，要把 RMT 外设通道 0 的输入信号 rmt\_sig\_in0（信号索引号 83）绑定到 GPIO40，请按照以下步骤操作（请注意 GPIO40 也叫做 MTDO 管脚）：
\begin{enumerate}
\item 置位 \hyperref[regdesc:GPIOFUNCNINSELCFGREG]{GPIO\_FUNC83\_IN\_SEL\_CFG\_REG} 寄存器的 \hyperref[fielddesc:GPIOSIGNINSEL]{GPIO\_SIG83\_IN\_SEL}，使能通过 GPIO 交换矩阵接收外部输入信号；
\item 将 \hyperref[regdesc:GPIOFUNCNINSELCFGREG]{GPIO\_FUNC83\_IN\_SEL\_CFG\_REG} 寄存器的 \hyperref[fielddesc:GPIOFUNCNINSEL]{GPIO\_FUNC83\_IN\_SEL} 字段设置为 40；
\item 置位 \hyperref[regdesc:IOMUXnREG]{IO\_MUX\_GPIO40\_REG} 寄存器的 \hyperref[fielddesc:IOMUXFUNIE]{IO\_MUX\_FUN\_IE}（使能输入模式）。
\end{enumerate}

{\bfseries 说明：}
\begin{itemize}
\item 同一个输入 pad 上可以同时绑定多个内部 input\_signals。
\item 置位 \hyperref[fielddesc:GPIOFUNCNININVSEL]{GPIO\_FUNC\regindex{y}\_IN\_INV\_SEL} 可以把输入的信号取反。
\item 无需将输入信号绑定到一个 pad 也可以使外设读取恒低或恒高电平的输入值。实现方式为选择特定的 \hyperref[fielddesc:GPIOFUNCNINSEL]{GPIO\_FUNC\regindex{y}\_IN\_SEL} 输入值而不是一个 GPIO 序号：
\begin {itemize}
\item 当 \hyperref[fielddesc:GPIOFUNCNINSEL]{GPIO\_FUNC\regindex{y}\_IN\_SEL} 是 \hexprefix{}3C 时， input\_signal\_\regindex{X} 始终为 0；
\item 当 \hyperref[fielddesc:GPIOFUNCNINSEL]{GPIO\_FUNC\regindex{y}\_IN\_SEL} 是 \hexprefix{}38 时， input\_signal\_\regindex{X} 始终为 1。
\end {itemize}
\end{itemize}

\subsection{简单 GPIO 输入}\label{IOMUXSimpleGPIOInput}

\hyperref[regdesc:GPIOINREG]{GPIO\_IN\_REG}/\hyperref[regdesc:GPIOIN1REG]{GPIO\_IN1\_REG} 寄存器存储着每一个 GPIO pad 的输入值。

任意 GPIO pad 的输入值都可以随时读取而无需为某一个外设信号配置 GPIO 交换矩阵。但是需要为 pad \regindex{X} 的 \hyperref[regdesc:IOMUXnREG]{IO\_MUX\_\regindex{n}\_REG} 寄存器配置 \hyperref[fielddesc:IOMUXFUNIE]{IO\_MUX\_FUN\_IE} 位以使能输入，如章节 \ref{sec:inputconfig} 所述。
\section{通过 GPIO 交换矩阵的外设输出}

\subsection{概述}

为实现通过 GPIO 交换矩阵输出外设信号，需要配置 GPIO 交换矩阵将输出索引号为 ( 0 \verb+~+ 11，14 \verb+~+ 18, ...) 的外设信号输出到 42 个 GPIO (0 \verb+~+ 21, 26 \verb+~+ 45)。具体输出索引号见交换矩阵表格 \ref{table:peripheral_signals}。

输出信号从外设输出到 GPIO 交换矩阵，然后到达 IO MUX。IO MUX 必须设置相应 pad 为 GPIO 功能，这样输出 GPIO 信号就能连接到相应 pad。


\begin{tiplisting}
输出索引号为 223 \verb+~+ 227 的外设信号，可配置为从一个 GPIO 管脚输入后，直接由另一个 GPIO 管脚输出。
\end{tiplisting}
\vspace{-2em}
\subsection{功能描述}\label{IOMUXPeripheralOutput}

如图 \ref{Figure:iomux_gpio_matrix_all} 所示，对于信号输出，182 个输出信号中的某一个信号通过 GPIO 交换矩阵到达 IO MUX， 然后连接到某个 pad。

输出外设信号 \regindex{Y} 到某一 GPIO pad \regindex{X} 的步骤为：

\begin{enumerate}
    \item 在 GPIO 交换矩阵里配置 GPIO pad \regindex{X} 的 \hyperref[regdesc:GPIOFUNCNOUTSELCFGREG]{GPIO\_FUNC\regindex{x}\_OUT\_SEL\_CFG\_REG} 寄存器和 \hyperref[regdesc:GPIOENABLEREG]{GPIO\_ENABLE\_REG}[\regindex{x}] 字段。推荐使用相应 W1TS 或 W1TC 寄存器来更新 \hyperref[regdesc:GPIOENABLEREG]{GPIO\_ENABLE\_REG} 寄存器中的值：
    \begin{itemize}
        \item 设置 \hyperref[regdesc:GPIOFUNCNOUTSELCFGREG]{GPIO\_FUNC\regindex{x}\_OUT\_SEL\_CFG\_REG} 寄存器的 \hyperref[fielddesc:GPIOFUNCNOUTSEL]{GPIO\_FUNC\regindex{x}\_OUT\_SEL} 字段为外设输出信号 \regindex{Y} 的索引号 (\regindex{Y})。
        \item 要将信号强制使能为输出模式，需要将 GPIO pad \regindex{X} 对应的 \hyperref[regdesc:GPIOFUNCNOUTSELCFGREG]{GPIO\_FUNC\regindex{x}\_OUT\_SEL\_CFG\_REG} 寄存器的 \hyperref[fielddesc:GPIOFUNCNOENSEL]{GPIO\_FUNC\regindex{x}\_OEN\_SEL} 字段置位；同时需要将 \hyperref[regdesc:GPIOENABLEW1TSREG]{GPIO\_ENABLE\_W1TS\_REG} 或 \hyperref[regdesc:GPIOENABLE1W1TSREG]{GPIO\_ENABLE1\_W1TS\_REG} 中的相应字段置位。或者，将 \hyperref[fielddesc:GPIOFUNCNOENSEL]{GPIO\_FUNC\regindex{x}\_OEN\_SEL} 清零，即选择采用外设的输出使能信号，此时输出使能信号由内部逻辑功能决定。详细信息见表 \ref{table:peripheral_signals} 中“输出信号的输出使能信号”一栏。
        \item 清零 \hyperref[regdesc:GPIOENABLEW1TCREG]{GPIO\_ENABLE\_W1TC\_REG} 或 \hyperref[regdesc:GPIOENABLE1W1TCREG]{GPIO\_ENABLE1\_W1TC\_REG} 中相应位可以关闭 GPIO pad 的输出。
    \end{itemize}
    \item 要选择以开漏方式输出，可以设置 GPIO pad \regindex{X} 的 \hyperref[regdesc:GPIOPINNREG]{GPIO\_PIN\regindex{x}\_REG} 寄存器中的 \hyperref[fielddesc:GPIOPINNPADDRIVER]{GPIO\_PIN\regindex{x}\_PAD\_DRIVER} 位。
    \item 配置 IO MUX 寄存器来选择 GPIO 交换矩阵。配置 GPIO pad \regindex{X} 的 \hyperref[regdesc:IOMUXnREG]{IO\_MUX\_\regindex{x}\_REG} 的过程如下：
    \begin{itemize}
        \item 设置功能字段 \hyperref[fielddesc:IOMUXMCUSEL]{IO\_MUX\_MCU\_SEL} 为 GPIO pad \regindex{X} 的 IO MUX 功能（所有管脚的 Function 1，数值为 1）。
        \item 设置 \hyperref[fielddesc:IOMUXFUNDRV]{IO\_MUX\_FUN\_DRV} 字段为特定的输出强度值 (0 \verb+~+ 3)，值越大，输出驱动能力越强：
        \begin{itemize}
            \item 0: \verb+~+5 mA
            \item 1: \verb+~+10 mA
            \item 2: \verb+~+20 mA（默认值）
            \item 3: \verb+~+40 mA
        \end{itemize}
        \item 在开漏模式下，通过置位/清零 \hyperref[fielddesc:IOMUXFUNWPU]{IO\_MUX\_FUN\_WPU} 和 \hyperref[fielddesc:IOMUXFUNWPD]{IO\_MUX\_FUN\_WPD} 使能或关闭上拉/下拉电阻器。
    \end{itemize}
\end{enumerate}

{\bfseries 说明：}
\begin{itemize}
\item 某一个外设的输出信号可以同时从多个 pad 输出。
\item GPIO46 不可用于输出信号。
\item 置位 \hyperref[fielddesc:GPIOFUNCNOUTINVSEL]{GPIO\_FUNC\regindex{n}\_OUT\_INV\_SEL} 可以把输出的信号取反。
\end{itemize}

\subsection{简单 GPIO 输出}

GPIO 交换矩阵也可以用于简单 GPIO 输出，具体配置如下：
\begin{itemize}
    \item 设置 GPIO 交换矩阵 \hyperref[fielddesc:GPIOFUNCNOUTSEL]{GPIO\_FUNC\regindex{n}\_OUT\_SEL} 寄存器为特定的外设索引值 256  (\hexprefix{}100)；
    \item 设置 \hyperref[regdesc:GPIOOUTREG]{GPIO\_OUT\_REG}[31:0] 或者 \hyperref[regdesc:GPIOOUT1REG]{GPIO\_OUT1\_REG}[21:0] 寄存器中某一位的值为期望 GPIO 输出的值。
\end{itemize}

{\bfseries 说明：}
\begin{itemize}
\item \hyperref[regdesc:GPIOOUTREG]{GPIO\_OUT\_REG}[0] \verb+~+ \hyperref[regdesc:GPIOOUTREG]{GPIO\_OUT\_REG}[31] 对应 GPIO0 \verb+~+ GPIO31，\hyperref[regdesc:GPIOOUTREG]{GPIO\_OUT\_REG}[25:22] 无效；
\item \hyperref[regdesc:GPIOOUT1REG]{GPIO\_OUT1\_REG}[0] \verb+~+ \hyperref[regdesc:GPIOOUT1REG]{GPIO\_OUT1\_REG}[13] 对应 GPIO32 \verb+~+ GPIO45，\hyperref[regdesc:GPIOOUT1REG]{GPIO\_OUT1\_REG}[21:14] 无效。
\item 推荐使用相应的 W1TS (write 1 to set) 和 W1TC (write 1 to clear) 寄存器，例如：\hyperref[fielddesc:GPIOOUTW1TS]{GPIO\_OUT\_W1TS}/\hyperref[fielddesc:GPIOOUTW1TC]{GPIO\_OUT\_W1TC} 来改写 \hyperref[regdesc:GPIOOUTREG]{GPIO\_OUT\_REG} 或 \hyperref[regdesc:GPIOOUT1REG]{GPIO\_OUT1\_REG} 中的值。
\end{itemize}

\subsection{Sigma Delta 调制输出 (SDM)}

\subsubsection{功能描述}
182 个数字外设输出中有 8 个通道支持 1 比特二阶 Sigma Delta 调制输出，这 8 个通道对应的外设输出信号索引号为 100 \verb+~+ 107，默认使能输出。该调制器实现输出可配占空比的 PDM（脉冲密度调制）信号。二阶 Sigma Delta 调制器传输函数为：

\[
    H(z)=\textrm{X(z)z$^{-1}$} + \textrm{E(z)(1-z$^{-1}$)$^{2}$}
    \]

E(z) 为量化误差，X(z) 为输入。\\
Sigma Delta 调制器内部支持对 APB\_CLK 的 1 \verb+~+ 256 倍分频：

\begin{itemize}
    \item 置位 \hyperref[fielddesc:GPIOSDFUNCTIONCLKEN]{GPIOSD\_FUNCTION\_CLK\_EN} 使能调制器时钟；
    \item 配置寄存器 \hyperref[fielddesc:GPIOSDSDNPRESCALE]{GPIOSD\_SD\regindex{n}\_PRESCALE} 实现分频。\regindex{n} 为 0 \verb+~+ 7，对应 8 个通道。
\end{itemize}
分频后的时钟周期为调制器输出单位脉冲的周期。

\hyperref[fielddesc:GPIOSDSDNIN]{GPIOSD\_SD\regindex{n}\_IN} 为有符号数，范围为 [-128, 127]，配置此寄存器控制输出 PDM 信号的占空比 $^{1}$。

\begin{itemize}
    \item \hyperref[fielddesc:GPIOSDSDNIN]{GPIOSD\_SD\regindex{n}\_IN} = -128，调制器输出信号占空比为 \begin{math}0\% \end{math}；
    \item \hyperref[fielddesc:GPIOSDSDNIN]{GPIOSD\_SD\regindex{n}\_IN} = 0，调制器输出信号占空比接近  \begin{math}50\% \end{math}；
    \item \hyperref[fielddesc:GPIOSDSDNIN]{GPIOSD\_SD\regindex{n}\_IN} = 127，调制器输出信号占空比接近  \begin{math}100\% \end{math}。
\end{itemize}

PDM 信号占空比计算公式为：

    \[
    Duty\_Cycle=\frac{\hyperref[fielddesc:GPIOSDSDNIN]{GPIOSD\_SD\regindex{n}\_IN} + 128} {256}
    \]

\begin{tiplisting}
对 PDM 信号来说，占空比是指在若干脉冲周期内（比如 256 个脉冲周期），高电平占整个统计周期的比值。
\end{tiplisting}

\subsubsection{配置方法}

SDM 的配置方法如下：

\begin{itemize}
    \item 将 SDM 输出经 GPIO 交换矩阵连接至 pad，见 \ref{IOMUXPeripheralOutput} 章节。
    \item 置位 \hyperref[fielddesc:GPIOSDFUNCTIONCLKEN]{GPIOSD\_FUNCTION\_CLK\_EN}，使能 SDM 时钟。
    \item 配置 \hyperref[fielddesc:GPIOSDSDNPRESCALE]{GPIOSD\_SD\regindex{n}\_PRESCALE} 寄存器设置时钟分频系数。
    \item 配置 \hyperref[fielddesc:GPIOSDSDNIN]{GPIOSD\_SD\regindex{n}\_IN} 设置 SDM 输出信号的占空比。
\end{itemize}

\section{专用 GPIO}\label{section:dedicatedGPIO}

\subsection{概述}
专用 GPIO 模块是为 CPU 与 GPIO 交换矩阵和 IO MUX 交互而设计的专用通道，该模块包括 8 个输入输出通道，其中输入通道对应的外设输入信号索引号为 235 \verb+~+ 242；输出通道对应的外设输出信号索引号也为 235 \verb+~+ 242，输出通道默认使能输出。

\begin{figure}[h!]
    \centering
    \includegraphics[width=1.0\textwidth]{\modulefiles/figures/dedicated_gpio}
    \caption{专用 GPIO 模块图}
    \label{Figure:dedicated_gpio}
\end{figure}

图 \ref{Figure:dedicated_gpio} 为专用 GPIO 结构图。用户可通过置位寄存器 \hyperref[regdesc:SYSTEMCPUPERICLKENREG]{SYSTEM\_CPU\_PERI\_CLK\_EN\_REG} 中的 \hyperref[fielddesc:SYSTEMCLKENDEDICATEDGPIO]{SYSTEM\_CLK\_EN\_DEDICATED\_GPIO} 位使能专用 GPIO 模块的时钟，也可以通过先置位，再清零寄存器 \hyperref[regdesc:SYSTEMCPUPERIRSTENREG]{SYSTEM\_CPU\_PERI\_RST\_EN\_REG} 中的 \hyperref[fielddesc:SYSTEMRSTENDEDICATEDGPIO]{SYSTEM\_RST\_EN\_DEDICATED\_GPIO} 位复位专用 GPIO 模块。详细信息可参考章节 \ref{mod:sysreg} \textit{\nameref{mod:sysreg}} 中
表 \ref{table:peripheralcontrollingbit} \textit{\nameref{table:peripheralcontrollingbit}}。

\subsection{主要特性}

专用 GPIO 模块具备以下特性：
\begin{itemize}
    \item 8 个输出通道和 8 个输入通道
    \item 每个通道均可通过寄存器或 CPU 指令访问
    \item 输入通道延迟时间可配置
    \item 输入通道支持中断
\end{itemize}

\subsection{功能描述}

用户可通过寄存器访问专用 GPIO，或直接调用特定的 CPU 指令访问专用 GPIO。

如图 \ref{Figure:dedicated_gpio} 所示，对于输出通道，专用 GPIO 模块提供两种方式来驱动 GPIO。用户可通过配置寄存器 \hyperref[regdesc:DEDICGPIOOUTCPUREG]{DEDIC\_GPIO\_OUT\_CPU\_REG} 中相应位选择具体的 GPIO 输出驱动方式：
\begin{itemize}
    \item \hyperref[fielddesc:DEDICGPIOOUTCPUSEL0]{DEDIC\_GPIO\_OUT\_CPU\_SEL\regindex{n}} = 0，则通过配置寄存器的方式来驱动 GPIO 的输出；
    \item \hyperref[fielddesc:DEDICGPIOOUTCPUSEL0]{DEDIC\_GPIO\_OUT\_CPU\_SEL\regindex{n}} = 1，则通过 CPU 指令来驱动 GPIO 的输出。
\end{itemize}
对于输入通道，专用 GPIO 模块也提供两种方式来获取专用 GPIO 输入值：
\begin{itemize}
    \item 通过寄存器查询专用 GPIO 输入值；
    \item 通过 CPU 指令读取专用 GPIO 的输入值。
\end{itemize}

\subsubsection{通过寄存器访问专用 GPIO}

寄存器可通过下面三种方式驱动 GPIO 输出：
\begin{itemize}
    \item 直接写 GPIO 输出值，通过配置寄存器 \hyperref[regdesc:DEDICGPIOOUTDRTREG]{DEDIC\_GPIO\_OUT\_DRT\_REG} 实现；
    \item MASK 写 GPIO 输出值，通过配置寄存器 \hyperref[regdesc:DEDICGPIOOUTMSKREG]{DEDIC\_GPIO\_OUT\_MSK\_REG} 实现；
    \item 单独 BIT 写 GPIO 输出值，通过配置寄存器 \hyperref[regdesc:DEDICGPIOOUTIDVREG]{DEDIC\_GPIO\_OUT\_IDV\_REG} 实现。
\end{itemize}

软件可通过寄存器 \hyperref[regdesc:DEDICGPIOOUTSCANREG]{DEDIC\_GPIO\_OUT\_SCAN\_REG} 查询 GPIO 的状态，即为图 \ref{Figure:dedicated_gpio} 中的 gpio\_out\_status。


软件可以通过读寄存器 \hyperref[regdesc:DEDICGPIOINSCANREG]{DEDIC\_GPIO\_IN\_SCAN\_REG} 获取专用 GPIO 输入值，即为图 \ref{Figure:dedicated_gpio} 中的 gpio\_in\_status。

专用 GPIO 模块支持对输入信号零延时及一个时钟/两个时钟/三个时钟周期的延时，可通过配置寄存器 \hyperref[regdesc:DEDICGPIOINDLYREG]{DEDIC\_GPIO\_IN\_DLY\_REG} 来控制各个通道的延时。 专用 GPIO 模块支持可配置的中断方式来指示 GPIO 的输入状态。配置寄存器 \hyperref[regdesc:DEDICGPIOINTRRCGNREG]{DEDIC\_GPIO\_INTR\_RCGN\_REG} 设置输入信号中断产生方式，包括：
\begin{itemize}
    \item 0/1：无中断
    \item 2：低电平触发
    \item 3：高电平触发
    \item 4：下降沿触发
    \item 5：上升沿触发
    \item 6/7：沿触发
\end{itemize}


\subsubsection{通过 CPU 指令访问专用 GPIO}\label{cpu_instructor}

CPU 也可通过指令来读/写专用 GPIO，共包括六条指令：
\begin{itemize}
    \item 置位输出通道中相应位\\
    汇编语法：SET\_BIT\_GPIO\_OUT mask  \\
    寻址方式：立即数寻址 \\
功能为写 1 置位用户寄存器 GPIO\_OUT 对应的位，mask 的位宽为 8。即 GPIO\_OUT 中某些位将被置位，这些位即为 mask 中为 1 的位，其余位则保持不变。GPIO\_OUT 对应图 \ref{Figure:dedicated_gpio} 中的 gpio\_out。
    \item 清除输出通道中相应位\\
    汇编语法：CLR\_BIT\_GPIO\_OUT mask \\
    寻址方式：立即数寻址 \\
功能为写 1 清零用户寄存器 GPIO\_OUT 对应的位，mask 的位宽为 8。即 GPIO\_OUT 中某些位将被清零，这些位即为 mask 中为 1 的位，其余位则保持不变。
    \item 以 MASK 方式置位或清除输出通道中相应位\\
    汇编语法：WR\_MASK\_GPIO\_OUT value, mask \\
    寻址方式：寄存器寻址 \\
功能为按 mask 方式向用户寄存器 GPIO\_OUT 写入 value。value 位宽为 8，表示将要写入的值；mask 位宽为 8，表示需要操作的 GPIO\_OUT。只有 mask 中为 1 的位被更新，即更新为 value 的值。如 mask \hexprefix{}03 (0000 0011), 表示写入值只对 GPIO\_OUT[0] 与 GPIO\_OUT[1] 有效。
    \item 写地址寄存器 art 的值到输出通道\\
    汇编语法：WUR.GPIO\_OUT art \\
    寻址方式：寄存器寻址 \\
功能为写地址寄存器 art 的值到用户寄存器 GPIO\_OUT。寄存器 art 的位宽为 32，使用该指令时只有低 8 位有效。
    \item 读输出通道的值到 arr 寄存器\\
    汇编语法：RUR.GPIO\_OUT arr \\
    寻址方式：寄存器寻址 \\
功能为读用户寄存器 GPIO\_OUT 的值到地址寄存器 arr。寄存器 arr 的位宽为 32，使用该指令时只有低 8 位有效。
    \item 读输入通道的值至地址寄存器 I \\
    汇编语法：GET\_GPIO\_IN I \\
    寻址方式：寄存器寻址 \\
功能为读用户寄存器 GPIO\_IN 的值到地址寄存器 I。寄存器 I 的位宽为 32, 其高 24 位为 0，低 8 位对应用户寄存器 GPIO\_IN 的值。
\end{itemize}

\section{IO MUX 的直接 I/O 功能}

\subsection{概述}

快速信号如 SPI、JTAG 等会旁路 GPIO 交换矩阵以实现更好的高频数字特性。所以高速信号会直接通过 IO MUX 输入和输出。

这样比使用 GPIO 交换矩阵的灵活度要低，即每个 GPIO pad 的 IO MUX 寄存器只有较少的功能选择，但可以实现更好的高频数字特性。

\subsection{功能描述}

对于外设输入信号，I/O 旁路 GPIO 交换矩阵必须配置两个寄存器：

\begin{enumerate}
\item GPIO pad 的 \hyperref[fielddesc:IOMUXMCUSEL]{IO\_MUX\_MCU\_SEL} 必须设置为相应的 pad 功能，章节 \ref{sec:iomux_pads} 列出了 pad 功能。
\item 对于输入信号，必须将 \hyperref[fielddesc:GPIOSIGNINSEL]{GPIO\_SIG\regindex{n}\_IN\_SEL} 寄存器设置为低电平，则输入信号直接连接到外设。
\end{enumerate}

对于外设输出信号，I/O 旁路 GPIO 交换矩阵只需将 GPIO pad 的 \hyperref[fielddesc:IOMUXMCUSEL]{IO\_MUX\_MCU\_SEL} 配置为相应的 pad 功能即可。

\begin{tiplisting}
并非所有外设输入/输出信号均可直接通过 IO MUX 连接到外设，某些输入/输出信号只能通过 GPIO 交换矩阵连接到外设。
\end{tiplisting}

\section{RTC IO MUX 的低功耗和模拟 I/O 功能}

\subsection{概述}

ESP32-S2 中有 22 个 GPIO 管脚具有低功耗（低功耗 RTC）性能和模拟功能，由 RTC 子系统控制。这些功能不使用 IO MUX 和 GPIO 交换矩阵，而是使用 RTC IO MUX 将 I/O 指向 RTC 子系统。

当这些管脚被配置为 RTC GPIO 管脚，作为输出管脚时仍然能够在芯片处于 Deep-sleep 模式下保持输出电平值或者作为输入管脚使用时可以将芯片从 Deep-sleep 中唤醒。

章节 \ref{sec:rtcmux_pins} 列出了 RTC\_MUX 管脚和功能。

\subsection{功能描述}

每个 pad 的模拟和 RTC 功能是由 \hyperref[regdesc:RTCIOTOUCHPADNREG]{RTCIO\_TOUCH\_PAD\regindex{n}\_REG} 寄存器中的 \hyperref[fielddesc:RTCIOTOUCHPADNMUXSEL]{RTCIO\_TOUCH\_PAD\regindex{n}\_MUX\_SEL} 位控制的。此位默认置为 0，通过 IO MUX 子系统输入输出信号，如前文所述。

如果置位 \hyperref[fielddesc:RTCIOTOUCHPADNMUXSEL]{RTCIO\_TOUCH\_PAD\regindex{n}\_MUX\_SEL}，则输入输出信号会经过 RTC 子系统。在这种模式下，\\\hyperref[regdesc:RTCIOTOUCHPADNREG]{RTCIO\_TOUCH\_PAD\regindex{n}\_REG} 寄存器用于数字 I/O，pad 的模拟功能也可以实现。章节 \ref{sec:rtcmux_pins} 列出了 RTC 管脚的功能。

表 \ref{sec:rtcmux_pins} 列出了 GPIO pad 与相应的 RTC 管脚和模拟功能的映射关系。请注意 \hyperref[regdesc:RTCIOTOUCHPADNREG]{RTCIO\_TOUCH\_PAD\regindex{n}\_REG} 寄存器使用的是 RTC GPIO 管脚的序号，不是 GPIO pad 的序号。

\section{Light-sleep 模式管脚功能}

当 ESP32-S2 处于 Light-sleep 模式时管脚可以有不同的功能。如果某一 GPIO pad 的 \hyperref[regdesc:IOMUXnREG]{IO\_MUX\_\regindex{n}\_REG} 寄存器中 \hyperref[fielddesc:IOMUXSLPSEL]{IO\_MUX\_SLP\_SEL} 位置为 1，芯片处于 Light-sleep 模式下将由另一组不同的寄存器控制 pad。

\subfile{\modulefiles/iomux_sleep_functions_cn.tex}

如果 \hyperref[fielddesc:IOMUXSLPSEL]{IO\_MUX\_SLP\_SEL} 置为 0，则芯片在正常工作和 Light-sleep 模式下，管脚的功能一样。

\begin{tiplisting}
正常工作模式或者 \hyperref[fielddesc:IOMUXSLPSEL]{IO\_MUX\_SLP\_SEL} = 0 时，输出使能配置请参考 \ref{IOMUXPeripheralOutput} 章节。

\end{tiplisting}

\section{Pad Hold 特性}\label{PadHold}

每个 IO pad（包括 RTC pad）都有单独的 hold 功能，由 RTC 寄存器控制。pad 的 hold 功能被置上后，pad 在置上 hold 那一刻的状态被强制保持，无论内部信号如何变化，修改 IO MUX 配置或者 GPIO 配置，都不会改变 pad 的状态。应用如果希望在看门狗超时触发内核复位时或 Deep-sleep 触发内核复位时 pad 的状态不被改变，就需要提前把 hold 置上。
\begin{tiplisting}
\vspace{-2em}
\begin{itemize}
\item 对于数字 pad 而言，若要在深度睡眠掉电之后保持 pad 输入输出的状态值，需要在掉电之前把寄存器 \hyperref[fielddesc:RTCCNTLDGPADFORCEUNHOLD]{RTC\_\-CNTL\_\-DG\_\-PAD\_\-FORCE\_\-UNHOLD} 设置成 0。对于 RTC pad 的输入输出值，由寄存器 \hyperref[regdesc:RTCCNTLPADHOLDREG]{RTC\_\-CNTL\_\-PAD\_\-HOLD\_\-REG} 中相应的位来控制 Hold 和 Unhold pad 的值。
\item 在芯片被唤醒之后，若要关闭 Hold 功能，将寄存器 \hyperref[fielddesc:RTCCNTLDGPADFORCEUNHOLD]{RTC\_CNTL\_DG\_PAD\_FORCE\_UNHOLD} 设置成 1。若想继续保持 pad 的值，可把 \hyperref[regdesc:RTCCNTLPADHOLDREG]{RTC\_CNTL\_PAD\_HOLD\_REG} 寄存器中相应的位设置成 1。

\end{itemize}
\end{tiplisting}

\vspace{-3em}
\section{I/O Pad 供电}\label{I/O Pad Power Supplies}

IO pad 供电请参考 \href{https://www.espressif.com/sites/default/files/documentation/esp32-s2_datasheet_cn.pdf}{《ESP32-S2 技术规格书》}中管脚定义章节。

\subsection{电源管理}

ESP32-S2 的数字管脚可分为如下 4 种不同的电源域。
\begin{itemize}
\item VDD3P3\_RTC\_IO：RTC 和 CPU 的输入电源
\item VDD3P3\_CPU：CPU 的输入电源
\item VDD3P3\_RTC：RTC 模拟的输入电源
\item VDD\_SPI：可配置为输入电源或输出电源
\end{itemize}

VDD\_SPI 可配置使用一个内置 LDO，该内置 LDO 的输入和输出均为 1.8 V。如未使能 LDO，VDD\_SPI 可以与 VDD3P3\_RTC\_IO 连接在相同的电源上。

VDD\_SPI 的具体配置由 GPIO45 的 Strapping 值决定，用户可通过 eFuse 或寄存器修改 VDD\_SPI 的配置。请参考\href{https://www.espressif.com/sites/default/files/documentation/esp32-s2_datasheet_cn.pdf}{《ESP32-S2 技术规格书》}中的电源管理章节和 Strapping 管脚章节查看更多信息。



\section{外设信号列表}
\label{sec:peripheral_signals}

表 \ref{table:peripheral_signals} 列出了 GPIO 交换矩阵的外设输入输出信号。

 \subfile{\modulefiles/peripheral_signals_cn.tex}

\section{IO MUX Pad 列表}
\label{sec:iomux_pads}
表 \ref{table:iomux_pads} 列出了每个 I/O pad 的 IO MUX 功能。

 \subfile{\modulefiles/iomux_pad_summary_cn.tex}

\textbf{复位配置}

“Reset”一栏是每个 pad 复位后的默认配置。

\begin{itemize}
\item \textbf{0} - IE=0（输入关闭）
\item \textbf{1} - IE=1（输入使能）
\item \textbf{2} - IE=1, WPD=1（输入使能，下拉电阻使能）
\item \textbf{3} - IE=1, WPU=1（输入使能，上拉电阻使能）
\end{itemize}

\textbf{说明}

\begin{itemize}
\item \textbf{R} - Pad 通过 RTC IO MUX 具有 RTC/模拟功能。
\item \textbf{I} - Pad 只能配置为输入 GPIO。
\end{itemize}

请参考 \href{https://www.espressif.com/sites/default/files/documentation/esp32-s2_datasheet_cn.pdf}{《ESP32-S2 技术规格书》}的附录查看管脚功能的完整表格。


\section{RTC IO MUX 管脚清单}
\label{sec:rtcmux_pins}

表 \ref{table:rtcmux pads} 列出了 RTC 管脚和对应 GPIO pad。

\subfile{\modulefiles/rtcmux_pad_summary_cn.tex}

\begin{tiplisting}
\vspace{-2em}
\begin{enumerate}
    \item TOUCH\_PAD0 为内部通道，其模拟功能未引出至外部管脚。
    \item 有关 sar\_i2c\_xx 的配置信息，请参考章节 \ref{mod:ulp} \textit{\nameref{mod:ulp}}：RTC I2C 控制器。
\end{enumerate}
\end{tiplisting}



\section{基地址}\label{label base addr}
用户可通过表 \ref{table:gprtcio_base_addr} 所示的基地址访问本章节所述模块。更多信息见章节 \ref{mod:sysmem} \textit{\nameref{mod:sysmem}}。

\begin{longtable}[c]{ |l|l|C{2.5cm}| }
\caption{模块基地址}\label{table:gprtcio_base_addr}\\
\hline
\rowcolor{lightgray}
    模块 &   访问总线      & 基地址 \\\hline
     \multirow{2}*{GPIO}   & PeriBUS1     & \hexprefix{}3F404000 \\\cline{2-3}
     & PeriBUS2       & \hexprefix{}60004000 \\\hline
     \multirow{2}*{IO MUX}   & PeriBUS1     & \hexprefix{}3F409000 \\\cline{2-3}
     & PeriBUS2       & \hexprefix{}60009000 \\\hline
     \multirow{1}*{GPIOSD} &  PeriBUS2  & \hexprefix{}60004F00 \\\hline
     \multirow{1}*{专用 GPIO} &  PeriBUS1  & \hexprefix{}3F4CF000 \\\hline
    \multirow{2}*{RTCIO}   & PeriBUS1     & \hexprefix{}3F408400 \\\cline{2-3}
     & PeriBUS2       & \hexprefix{}60008400 \\\hline
\end{longtable}

\hypertarget{iomuxgpio-reg-summ}{}
\section{寄存器列表}


\subsection{GPIO 交换矩阵寄存器列表}

请注意，下面的地址是相对于 GPIO 基地址的地址偏移量（相对地址），请参阅章节 \ref{label base addr} 获取有关 GPIO 基地址的信息。
\label{section:4.12}

\subfile{\modulefiles/gpio_reg_summary_cn.tex}

\subsection{IO MUX 寄存器列表}

请注意，下面的地址是相对于 IO MUX 基地址的地址偏移量（相对地址），请参阅章节 \ref{label base addr} 获取有关 IO MUX 基地址的信息。

\subfile{\modulefiles/io_mux_reg_summary_cn.tex}

\subsection{SDM 寄存器列表}

请注意，下面的地址是相对于 GPIOSD 基地址的地址偏移量（相对地址），请参阅章节 \ref{label base addr} 获取有关 GPIOSD 基地址的信息。

\subfile{\modulefiles/gpio_sd_reg_summary_cn.tex}

\subsection{专用 GPIO 寄存器列表}

请注意，下面的地址是相对于专用 GPIO 基地址的地址偏移量（相对地址），请参阅章节 \ref{label base addr} 获取有关专用 GPIO 基地址的信息。

\subfile{\modulefiles/dedicated_gpio_summary_cn.tex}

\subsection{RTC IO MUX 寄存器列表}
请注意，下面的地址是相对于 RTCIO 基地址的地址偏移量（相对地址），请参阅章节 \ref{label base addr} 获取有关 RTCIO 基地址的信息。
\subfile{\modulefiles/rtc_io_reg_summary_cn.tex}

\section{寄存器}


\subsection{GPIO 交换矩阵寄存器}
\subfile{\modulefiles/gpio_reg_cn.tex}

\subsection{IO MUX 寄存器}
\subfile{\modulefiles/io_mux_reg_cn.tex}

\subsection{SDM 寄存器}
\subfile{\modulefiles/gpio_sd_reg_cn.tex}

\subsection{专用 GPIO 寄存器}
\subfile{\modulefiles/dedicated_gpio_reg_cn.tex}

\subsection{RTC IO MUX 寄存器}
\subfile{\modulefiles/rtc_io_reg_cn.tex}

\end{document}
